
% Variant "tikzmark"
%    - Stores in .aux-File
%----------------------------------------------------
%====================================================
%####################################################
\documentclass[a4paper]{article}
\usepackage{lipsum,showframe}
\usepackage{tikz}
\usetikzlibrary{tikzmark}

\newlength{\lastpagegap}% reserve global name

\makeatletter
\newcommand{\pagespace}{%
  \hrule height0pt %terminate last paragraph, is needed
  \strut\tikzmark{lastpage}
  \vfill
  \noindent\begin{tikzpicture}[remember picture,overlay]
    \coordinate (lastpage) at (pic cs:lastpage);
    \pgfextracty{\lastpagegap}{\pgfpointanchor{lastpage}{center}}
    \global\advance\lastpagegap by \ht\strutbox
    %\node[above right]{\the\lastpagegap};
  \end{tikzpicture}%
  \write\@auxout{\string\global\string\lastpagegap=\the\lastpagegap}% store in aux file
 }
\makeatother

\begin{document}
\the\lastpagegap

\lipsum

Random text\footnote{A random footnote.}.

\begin{figure}[t] \rule{3cm}{3cm} \caption{A small black square} \end{figure}

\pagespace\smash{\rule{10pt}{\lastpagegap}}
\end{document}
%####################################################
%====================================================
%----------------------------------------------------












% Variant "Calculation"
%    - Creates Cmd in .aux-File to then store in own auxiliary-File
%----------------------------------------------------
%====================================================
%####################################################
\documentclass[a4paper]{article}
\usepackage{showframe}
\usepackage{lipsum}

\makeatletter
% write information about free space to \jobname.xyz
\AtEndDocument{\if@filesw\newwrite\tf@xyz
    %                                 ^^^
    \immediate\openout\tf@xyz\jobname.xyz\fi
    %                     ^^^         ^^^
}
\newdimen\remainingheight
\newdimen\minimumheight
\newif\ifendonnewpage
\endonnewpagefalse
\newcommand*{\calcremaining}{%
    \par
    % in case a new page just started report a remaining height of 0pt
    \ifdim\pagegoal=\maxdimen
        \remainingheight\z@
        \endonnewpagetrue
    \else
        \remainingheight\pagegoal
        \advance\remainingheight by -\pagetotal
        \advance\remainingheight by -\lineskip
        \endonnewpagefalse
    \fi
}
\newcommand*{\getremainingheight}{%
    \calcremaining
    % if not enough space, report 0pt remaining space
    \ifdim\remainingheight<\minimumheight
        \remainingheight\z@
    \else
        % important, otherwise the remaining height reported would be slightly too large
        \noindent
        \calcremaining
    \fi
    % write to own aux file via main aux
    \immediate\write\@auxout{\string\@writefile{xyz}{\string\remainingheight\the\remainingheight}}%
    %                                           ^^^
    \immediate\write\@auxout{\string\@writefile{xyz}{\ifendonnewpage\string\endonnewpagetrue\else\string\endonnewpagefalse\fi}}%
    %                                           ^^^
}
% just load the file to show the values in this document
% this will of course need two runs
\AtBeginDocument{\InputIfFileExists{\jobname.xyz}{}{}}
%                                            ^^^
\AtEndDocument{\getremainingheight
% just for testing, will of course produce a third page, if the
% remaining height is less then \baselineskip
%    \ifdim\remainingheight>0pt\rule{10pt}{\remainingheight}\fi
}
\makeatother

% set minimum space to be reported, less then that will be reported as 0pt
\setlength{\minimumheight}{\baselineskip}

\begin{document}
\lipsum

Remaining height: \the\remainingheight; ended on new page: \ifendonnewpage yes\else no\fi

Random text\footnote{A random footnote.}.

% see what happens if page is full
%\lipsum[4]
%\lipsum[4]
%One more line to fill the page.

\begin{figure}[t] \rule{2cm}{2cm} \caption{A small black square} \end{figure}

\begin{figure}[b] \rule{3cm}{3cm} \caption{A black square} \end{figure}

% this will produce a float page, which is not recognized at all
%\begin{figure}[p] \rule{4cm}{4cm} \caption{A big black square} \end{figure}

% see what happens, if there is a newpage at the end
%\newpage
\end{document}
%####################################################
%====================================================
%----------------------------------------------------